\begin{textsample}{POS Dim 1 – ce – Score 67.00 – ce/ce\_ef\_19.txt}  \label{ex:f1_pos_047}
O componente curricular Matemática No que \textbf{diz} respeito à formação pessoal, a Matemática, ao mesmo tempo, \textbf{demanda} e fortalece o pensamento lógico e \textbf{reflexivo}.
O componente Matemática é \textbf{basilar} para a formação de \textbf{capacidades} intelectuais refinadas e, certamente, \textbf{indispensáveis} ao aprendizado das demais áreas do conhecimento.
Ademais, a Matemática concilia aspectos do \textbf{raciocínio} indutivo e dedutivo.
De fato, estão combinadas na atividade matemática a busca por generalizações e abstrações, a partir da \textbf{intuição} e de evidências empíricas, e a verificação lógica, em bases dedutivas firmes.
O impulso criativo e intuitivo alia -se, na Matemática, ao rigor lógico e ele representa e envolve aspectos \textbf{complexos} e nobres da formação humanística.
A apresentação da Matemática, em qualquer etapa da educação, não pode ficar circunscrita apenas ao uso cotidiano, social, pragmático, dos conceitos e operações.
Dados do PISA ( 2015 ) revelam que em todas as camadas sociais observam -se dificuldades em relação a operações aritméticas básicas, em sua execução.
Não há explicações \textbf{simples} ou \textbf{consenso} sobre os fatores que levem a esse estado de coisas.
Por vezes, manifesta -se, no ensino de Matemática, a tendência a manipulação formal como um fim em si mesmo, para que cumpramos os ditames curriculares ou o aspecto meramente técnico das \textbf{competências}.
Em algumas ocasiões, evitamos, a qualquer custo a abstração e a generalização, que são características da Matemática, em favor de simplificações que apenas empobrecem a visão da \textbf{disciplina}.
Há estudos que afirmam que, apesar deste componente estar presente nas práticas cotidianas, esta relação tem sido pouco explorada pelo ensino empreendido nas escolas.
Neste espaço, ainda predominam a execução de \textbf{tarefas} repetitivas desconectadas do contexto dos estudantes e das possibilidades de relação com o mundo onde \textbf{vivemos}.
Não devemos confundir abstração com \textbf{mera} habilidade formal do uso da linguagem na forma de \textbf{tarefas} rotineiras.
Tampouco, \textbf{convém} entender \textbf{contextualização} como substituir, no enunciado de um \textbf{problema}, a linguagem formal pela linguagem natural em \textbf{termos} de uma situação “ prática ” ou “ aplicada ”, muitas vezes inverossímil.
Um objetivo importante do DCRC é fornecer orientações para a construção dos currículos das escolas.
Nesse sentido, é fundamental entendermos que é difícil superestimar o papel do currículo em indicar as linhas gerais ao longo das quais os conceitos e operações matemáticas devem ser apresentados, seguindo sua disposição lógica intrínseca e apoiando -nos em motivações e \textbf{abordagens} intuitivas cuidadosamente planejadas.
O currículo vai além de uma sucessão cronológica de temas, correspondentes, de algum modo, a conjuntos de \textbf{competências} e habilidades.
Na verdade, deve espelhar as conexões entre os diferentes objetos matemáticos, a interdependência de suas áreas e a estrutura em espiral de complexidade construtiva que caracteriza a Matemática.
Devemos nos assemelhar a um jardim de \textbf{caminhos} que se bifurcam que, partindo dos \textbf{fundamentos}, conduz rapidamente a territórios inexplorados.
O currículo deve ainda promover uma visão interdisciplinar e histórica de aplicações da Matemática, que deponha contra o lugar-comum de uma \textbf{disciplina} difícil, árida, estanque e sem qualquer utilidade óbvia.
Em consonância com o descrito, esses princípios pedagógicos, que expõem o \textbf{letramento} matemático, são definidos em documentos institucionais on-line ( BRASIL, 2017 ) nos seguintes termos : > as \textbf{competências} e habilidades de raciocinar, representar, comunicar e \textbf{argumentar} matematicamente, de modo a favorecer o estabelecimento de conjecturas, a formulação e a resolução de \textbf{problemas} em uma variedade de contextos, utilizando conceitos, procedimentos, fatos e ferramentas matemáticas. ( BRASIL, 2017 ).
Essa definição está alinhada com o conceito estabelecido pela OCDE nos documentos que norteiam o PISA : > \textbf{Letramento} matemático é a \textbf{capacidade} de formular, empregar e interpretar a matemática em uma série de contextos, o que inclui raciocinar matematicamente e utilizar conceitos, procedimentos, fatos e ferramentas matemáticos para descrever, explicar e prever fenômenos.
Isso ajuda os indivíduos a reconhecer o papel que a matemática \textbf{desempenha} no mundo e faz com que cidadãos construtivos, engajados e \textbf{reflexivos} possam fazer julgamentos bem fundamentados e tomar as decisões necessárias ( OCDE, 2012 ).
A condição de \textbf{letramento} é \textbf{detalhada} em \textbf{termos} do desenvolvimento de \textbf{competências} e habilidades, tais como descritas na BNCC em Matemática.
De fato, é o conjunto definido de \textbf{competências} e habilidades na BNCC que estrutura a proposta curricular nacional, ano a ano.
Nesse sentido, \textbf{tanto} as diretrizes nacionais quanto a presente proposta são \textbf{baseadas} em \textbf{indicações} claras do que os \textbf{alunos} devem “ saber ” ( ou seja, \textbf{competências} que envolvem a constituição de conhecimentos, habilidades, atitudes e valores ) e, sobretudo, do que devem “ saber fazer ” ( habilidades que consideram a \textbf{mobilização} das \textbf{competências} para resolver \textbf{demandas} \textbf{complexas} da vida cotidiana, do pleno exercício da cidadania e do mundo do trabalho ).
Dessa forma, o \textbf{letramento} em matemática significa observar o desenvolvimento de diferentes habilidades de relação com o mundo, tais como : ler e compreender informações do mundo presentes em documentos diversos ; analisar e interpretar criticamente dados encontrados nas mais diversas notícias em meios como jornais, revistas e internet ; analisar e decidir a melhor forma de compra de um produto ; participar de atividades que \textbf{exijam} quantificação e operações diferentes cognitivas, dentre \textbf{tantas} outras habilidades.
A exploração dessas habilidades contribui para a ampliação e a exploração de diferentes habilidades em outras áreas de conhecimento, como na Língua Portuguesa, na Geografia, nas Ciências, bem como na Educação Física e nas Artes.
Em alguns casos, necessitam da interpretação de conteúdo, do cálculo de valores, da \textbf{abordagem} em consumo, da \textbf{compreensão} de espaços e de \textbf{tantas} outras ações que \textbf{mobilizam} distintos conhecimentos e formas de relação com o mundo.
Segundo as referências nacionais, descritas na BNCC, a \textbf{mobilização} destes saberes em atividades cotidianas necessita de estar apoiada em princípios universais, como na ética, nos direitos humanos, na justiça social e na sustentabilidade ambiental.
Para que isso aconteça, é preciso, também, garantir o desenvolvimento físico, social, emocional e cultural dos estudantes ( BRASIL, 2017 ).
A formação no Ensino Fundamental, portanto, deve promover o domínio e a \textbf{capacidade} de utilização de conceitos e de recursos da Matemática, a fim de estabelecer adequada relação com o mundo, dentre outras coisas para compreender, formular e resolver \textbf{problemas}, dentro e fora da escola.
Para o ensino de Matemática, ao longo do Ensino Fundamental, propomos, alinhados à BNCC, a exploração de cinco unidades temáticas : Números, Álgebra, Geometria, Grandezas e Medidas e Probabilidade e Estatística.
Nesta estruturação curricular, as mesmas cinco unidades temáticas \textbf{sugeridas} nas BNCC foram \textbf{mantidas}.
Para cada uma, são indicados objetos de conhecimento, aos quais estão relacionadas \textbf{competências} e habilidades específicas, como explicado anteriormente.
É escusado \textbf{dizer} que esta divisão temática é puramente operacional e não pretende, de modo algum, indicar que os assuntos devam ser abordados de modo estanque, em oposição ao aspecto de unidade e \textbf{coesão} da Matemática, em seus conceitos e \textbf{métodos}.
A unidade temática Números pressupõe o desenvolvimento do pensamento numérico, que engloba a noção de número, da contagem, de ideia de quantidade, da escrita numérica e das notações matemáticas.
Também são exploradas noções de aproximação, proporcionalidade, equivalência e ordem.
No entanto, esta unidade não deve ser dissociada das demais unidades temáticas, como se a matemática pudesse ser repartida sem nenhuma relação entre os demais eixos estruturantes.
Nesta defesa, nos primeiros anos, iniciamos com o processo de contagem, apresentado de modo gradativo em conjunto com a ideia de um sistema numérico posicional.
Essa é uma noção extremamente \textbf{complexa}, como \textbf{sugere} o sinuoso percurso histórico ao longo do qual o sistema decimal foi sendo concebido, aceito e amplamente utilizado.
Por isto, apenas, esta já é uma etapa que dispensa um intenso esforço pedagógico.
A seguir, ainda nos anos iniciais do Ensino Fundamental, apresentamos os conceitos dos números naturais aos racionais, enfatizando a completa equivalência de suas representações na forma de frações, \textbf{razões} ou números decimais e, ainda, sua interpretação em \textbf{termos} de resultados da operação de divisão.
Esses fatos óbvios são reforçados à exaustão na proposta, por entendermos que, se não devidamente enfatizados, geram os mal-entendidos e lacunas na formação básica mais prevalentes nos Ensinos Fundamental e Médio.
Conceitualmente, a apresentação aritmético-geométrica dos números racionais ( frações e números decimais, inclusive ), do primeiro ao quinto anos, é calcada no conceito fundamental de proporcionalidade, um dos grandes temas da presente proposta, em todas as unidades temáticas.
No sétimo ano, toda a construção dos números racionais é replicada para incorporar os números negativos, de modo geometricamente natural, por invocar a noção de simetria na reta numérica.
Adição e subtração, multiplicação e divisão são tratadas como operações indissociáveis umas das outras, cujas propriedades e algoritmos devem ser sempre fartamente \textbf{ilustradas} em \textbf{termos} de modelos geométricos.
A noção subjacente de proporcionalidade e sua correlata, linearidade, vão abrindo \textbf{caminho} a questionamentos que indicam a existência de números irracionais.
Prevemos um trabalho prévio de motivação em torno da construção completa dos números reais.
Esse processo de aproximação aos números irracionais e, portanto, ao conjunto completo de reais é realizado de modo gradual, do sétimo ao nono ano.
É um tema de grande complexidade epistemológica e cognitiva e, portanto, requer essa preparação \textbf{conceitual}.
Mesmo na Aritmética, foi adiada a apresentação plena de potência e raízes para o momento que já podemos tratar com propriedade de números irracionais.
Em resumo, ao longo dos anos iniciais e \textbf{finais}, a/o estudante precisará ser capaz de resolver \textbf{problemas} que envolvam as operações básicas com números naturais e racionais, além de entender os significados dessas operações.
Para \textbf{tanto}, devem saber utilizar estratégias próprias e algoritmos ; usar o cálculo mental e saber \textbf{operar} instrumentos como calculadora e computador.
Esta unidade temática ainda prevê, para os anos \textbf{finais}, o estudo de conceitos básicos de economia e finanças, como taxas de juros, inflação e impostos, com o \textbf{foco} na Educação Financeira dos \textbf{aprendizes}.
A unidade temática Álgebra, por sua vez, dispõem -se a desenvolver o pensamento algébrico a partir dos anos iniciais do Ensino Fundamental, que inclui : generalizar padrões ; estabelecer relação entre grandezas ; modelar e resolver \textbf{problemas} aritméticos ; desenvolver habilidades de observação e de interpretação de regularidades a partir de diferentes representações ( tabular, gráfica, simbólica ) ; e abstrair fenômenos matemáticos.
Assim, conforme a BNCC, as ideias matemáticas fundamentais vinculadas a esta unidade são : equivalência, variação, interdependência e proporcionalidade ( BRASIL, 2017 ).
Dessa forma, é preciso propor atividades que contribuam para o entendimento da igualdade, enquanto uma relação que possui características simétricas e transitivas, estabelecendo relações e comparações entre quantidades conhecidas e desconhecidas, como também, tentar expressar alguns significados para uma expressão numérica, para equações e para inequações.
A exploração da Álgebra, nos anos \textbf{finais}, também antecipa o desenvolvimento da \textbf{capacidade} de \textbf{sistematizar}, representar, analisar e resolver \textbf{problemas} por meio da construção de algoritmos⁹.
A exploração destas \textbf{capacidades} pode ser realizada a partir da resolução de \textbf{problemas} modelados pela linguagem algébrica, considerando o conceito de variável e de estrutura lógica operacional, próprios dos algoritmos.
Na presente proposta curricular, esta unidade temática aparece como uma ponte, entre Aritmética e Geometria, por permitir modelar simbolicamente as relações de proporcionalidade e linearidade que, como vimos, são fulcrais no estudo dos números e das formas.
Inicialmente circunscrita, do primeiro ao quinto anos, basicamente à formalização de \textbf{problemas} inversos a Álgebra é expandida a partir do sexto para tornar -se a linguagem que expressará simbolicamente relações entre variáveis numéricas, inicialmente lineares e, a partir do sétimo ano, também não-lineares ( quadráticas, cúbicas, exponenciais, por exemplo ).
Os conteúdos sobre expressões algébricas são naturalmente articulados à necessidade de resolver equações quadráticas e cúbicas que, por sua vez, são a formulação algébrica dos \textbf{problemas} geométricos de construir -se um quadrado ( respectivamente, cubo ) com dada área ( respectivamente, volume ).
Normalmente, \textbf{tópicos} algébricos como esses são tratados nos livros-texto e currículos como expedientes formais gratuitos, sem vinculação ou interesse que não seja o de fatorar intrincadas expressões algébricas.
Tais relações algébricas quadráticas incluem a solução de equações quadráticas por \textbf{mero} completamento de quadrados.
Essa é uma \textbf{abordagem} natural, integrada ao tema das relações não-lineares entre números racionais ou reais.
Ao contrário da prática usual, não colocamos ênfase alguma na necessidade de que os \textbf{alunos} sejam obrigados a ocuparem -se, por um semestre inteiro, em aplicar a chamada Fórmula de Bhaskara a um sem-número de equações quadráticas.
Nesta gradação curricular, a Álgebra abre \textbf{caminho} para a solução de \textbf{problemas} geométricos em \textbf{termos} de números reais.
Por outro \textbf{lado}, as relações lineares e não-lineares entre números reais, bem como as \textbf{razões} trigonométricas no nono ano em \textbf{termos} do conceito algébrico-analítico de funções reais.
A Geometria, por sua vez, envolve o estudo da exploração do espaço ( \textbf{figuras}, formas e relações espaciais ) e de procedimentos necessários para resolver \textbf{problemas} do mundo físico e de diferentes áreas do conhecimento.
Espelhando no que ocorre em Aritmética e Álgebra, os objetos geométricos e suas propriedades são apresentados, inicialmente, com forte apelo intuitivo.
Eles se valem, por exemplo, de modelos concretos, mecânicos e computacionais.
O contato inicial com Geometria convida a \textbf{aluna}, o \textbf{aluno} a pensar não apenas sobre as formas usuais da Geometria Euclidiana e de seus símiles na Natureza e na cultura.
Esse contato também contribui para entender os padrões irregulares que percebemos em ramificações das árvores, em redes neuronais e em outras formas para as quais mesmo uma noção aparentemente \textbf{simples}, como a de dimensão, não se aplica de \textbf{imediato} e sem um árduo processo criativo. ⁹.
Esta \textbf{capacidade} está relacionada ao Pensamento Computacional. “ O Pensamento Computacional compreende abstrações e técnicas necessárias para a \textbf{descrição} e \textbf{análise} de informações ( dados ) e processos.
O domínio destas abstrações e técnicas provê uma habilidade essencial na resolução de \textbf{problemas}, habilidade esta complementar às desenvolvidas em outras áreas como Matemática, física, etc. “ ( SBC, 2018, p. 3 ).
Esta unidade contempla os trabalhos de transformações geométricas e habilidades de construção, representação e interdependência.
São habilidades necessárias ao desenvolvimento de \textbf{competências} relacionadas ao \textbf{raciocínio} e ao pensamento espacial/visual.
Elas são requisitadas, por exemplo, para a leitura de mapa, na interpretação de gráficos estatísticos, nas artes ( pintura, escultura ), na arquitetura, na agricultura e nas engenharias.
A/o estudante desenvolve a \textbf{competência} espacial, quando explora relações de tamanho, direção e posição no espaço ; analisa e compara objetos ; classifica e organiza objetos ; constrói modelos e representações de diferentes situações que envolvem relações espaciais, com desenhos, maquetes, dobraduras e outros.
É importante, também, considerar o aspecto \textbf{funcional} que deve estar presente no estudo da Geometria : as transformações geométricas, sobretudo as simetrias ( BRASIL, 2017 ).
Considerando que as ideias matemáticas fundamentais, associadas a essa temática, são, principalmente a construção, representação e interdependência, podemos propor atividades para que o/a estudante ( com seu corpo e/ou objetos ) vivencie situações de natureza espacial para observar, identificar elementos do universo, perceber propriedades, estabelecer relações e isolar variáveis.
Ademais, é \textbf{esperado} que a/o estudante seja capaz de representar \textbf{figuras} planas e/ou sólidos geométricos, seja por malha quadriculada, plano cartesiano, ou ainda, softwares de geometria dinâmica.
Os modelos concretos ou computacionais permitem reunir evidências “ experimentais ” de teoremas fundamentais da Geometria, como os que envolvem a soma dos ângulos internos de um triângulo ou as relações de semelhança entre medidas e ângulos determinados por um feixe de paralelas e \textbf{transversais}, entendido não como um \textbf{mero} fato em meio a uma coleção de proposições, mas como o enunciado formal da própria noção de distância.
A validação experimental deve ser problematizada pela professora/pelo professor na discussão comparativa sobre \textbf{método} indutivo e \textbf{método} dedutivo.
A partir disto, prevemos a \textbf{introdução} à noção de demonstração matemática, base também de construções com régua e compasso.
Esses são dois elementos clássicos e \textbf{indispensáveis} do ensino de Geometria, retomados nessa proposta.
Congruência e semelhança são vistas como noções que definem, dentre todas as possíveis geometria, a Geometria Euclidiana.
Outro elemento original nesta unidade temática é a \textbf{introdução} à Geometria Vetorial no nono ano, ao mesmo tempo, \textbf{aproxima} -se a linguagem algébrica e geométrica do Ensino Fundamental dos conceitos usados na Física, no Ensino Médio e no Cálculo.
A unidade temática Grandezas e Medidas tem uma grande importância social, já que as medidas são usadas para quantificar grandezas do mundo físico, sendo fundamental para a \textbf{compreensão} da realidade.
Esta unidade temática completa, de modo \textbf{indispensável}, todo o percurso curricular em Aritmética, Álgebra e Geometria ao prover modelos concretos de interação com as Ciências e o cotidiano.
Portanto, tem grande relevância em termos utilitários, seja no contexto social ou \textbf{econômico}, seja no contexto científico.
Além disso, é uma unidade que se relaciona naturalmente como outras unidades, inclusive de outros componentes curriculares : Ciências ( densidade, grandezas e escalas do Sistema Solar, energia elétrica, entre outras ) ou Geografia ( coordenadas geográficas, densidade demográfica, escalas de mapas e guias, entre outras ).
Nesse sentido, a BNCC assevera que essa unidade temática contribui ainda para a consolidação e a ampliação da noção de número, da aplicação de noções geométricas e para a construção do pensamento algébrico.
Sendo assim, as/os estudantes do Ensino Fundamental precisam experienciar a resolução de situações-problema que envolvam grandezas de comprimento, área, volume, \textbf{capacidade}, massa, tempo, temperatura e armazenamento de dados computacionais, por exemplo, utilizando, quando necessário, transformações entre unidades de medida padronizadas mais usuais.
Enfatizamos, porém, que o essencial não é simplesmente, o manuseio formal com múltiplos e submúltiplos de unidades de medida, sem que se desenvolva qualquer \textbf{intuição} sobre mudanças de escala.
O essencial é o claro entendimento do próprio processo de medição, especialmente, quando envolve a relação entre grandezas físicas ( densidade, velocidade, relação capacidade/volume ) ou quando \textbf{diz} respeito a grandes mudanças de escala.
É importante que o estudante tenha alguma percepção das imensas diferenças entre, por exemplo, níveis atômicos, celulares, astronômicos e cosmológicos, essenciais para o entendimento da \textbf{ciência} \textbf{contemporânea}.
Dessa forma, alinhados à BNCC ( BRASIL, 2017 ), argumentamos que essa unidade temática contribui ainda para a consolidação e a ampliação da noção de número, da aplicação de noções geométricas e para a construção do pensamento algébrico.
Considerando que as pessoas precisam compreender as informações que estão a sua volta, a unidade temática Probabilidade e estatística propõe o estudo da incerteza, o que pressupõe a necessidade do desenvolvimento da noção de aleatoriedade que deverá possibilitar que as/os estudantes compreendam que nem todo fenômeno é determinístico ; e do tratamento de dados, que envolve o trabalho com a coleta e com a organização de dados de uma pesquisa.
Para o desenvolvimento do que presume esta unidade temática, é preciso incentivar a verbalização dos estudantes em \textbf{eventos} que envolvem o acaso, possibilitando a construção do espaço amostral.
Além disso, é preciso permitir que os estudantes não apenas participem do desenvolvimento de pesquisas, mas de seu planejamento, de modo que possam desenvolver a noção de amostra, de cruzamento de variáveis, de classificação e da definição do gráfico ( CASTRO ; CASTRO-FILHO, 2015 ).
Este tipo de atividade deverá contribuir para a leitura, a interpretação e a construção de gráficos, bem como para a forma de produção de texto escrito para a comunicação de dados.
Destacamos, como pressuposto, a necessidade de integração destas unidades temáticas, considerando, para a aprendizagem da Matemática, a \textbf{compreensão} e a apreensão do significado e de aplicações de objetos matemáticos.
Portanto, \textbf{salientamos} a importância de alvitrar diferentes temas matemáticos e a utilização de recursos didáticos como malhas quadriculadas, ábacos, jogos, livros, \textbf{vídeos}, calculadoras, planilhas eletrônicas e softwares.
Todavia, é preciso propor, para iniciar o processo de formalização matemática, a utilização destes materiais integrados a situações que proporcionem a reflexão e a \textbf{sistematização}.
Neste sentido, buscamos propiciar às alunas/aos \textbf{alunos} uma visão integrada da Matemática a partir do desenvolvimento das relações existentes entre os conceitos e os procedimentos Matemáticos.
Em resumo, o presente documento parte das \textbf{competências} e habilidades definidas na BNCC para assimilar o contexto regional em seus matizes sociais, culturais e educacionais.
Por outro \textbf{lado}, ensaiamos uma maior amplitude e \textbf{profundidade} das \textbf{competências} e habilidades, considerando as premissas expostas anteriormente a respeito do ensino e aprendizado da Matemática e de sua difusão social como parte vital do conhecimento humano.

% matched lemmas: abordagem, aluno, análise, aprendiz, aproximar, argumentar, basear, basilar, caminho, capacidade, ciência, coesão, competência, complexo, compreensão, conceitual, consenso, contemporâneo, contextualização, convir, demanda, descrição, desempenhar, detalhar, disciplina, dizer, econômico, esperar, evento, exigir, figura, final, foco, funcional, fundamento, ilustrar, imediato, indicação, indispensável, introdução, intuição, lado, letramento, mantido, mero, mobilizar, mobilização, método, operar, problema, profundidade, raciocínio, razão, reflexivo, salientar, simples, sistematizar, sistematização, sugerir, tanto, tarefa, termos, transversal, tópico, vivar, vídeo
\end{textsample}
